\ifdefined\gkver\else\def\gkver{student}\fi
\documentclass[\gkver]{gaokaozhenti}
\newcommand{\ccnum}[1]{{\fontspec{Noto Sans CJK JP}#1}}
\begin{document}

\chapter{2019年全国理综Ⅲ}
%% paperType: 理综物理部分

\begin{enumerate}
\item
%% number: 14
%% paperName: 2019年全国理综Ⅲ第 14 题 6 分
%% typeId: 1
%% type: 单选题
%% chapter: 电磁感应
%% point: 楞次定律
%% method: 无
%% score: 6
%% degree: 950
%% duplicateId: 0
%% body:
楞次定律是下列哪个定律在电磁感应现象中的具体体现？\xzanswer{D}
\fourchoices[answer=D]
{电阻定律}
{库仑定律}
{欧姆定律}
{能量守恒定律}

%% answer: D
%% memo:
\memoanswer{
楞次定律是电磁感应现象中判断感应电流方向的定律，不论它以何种方式表述，都是感应电流的磁场与原磁场之间的相互阻碍关系，而实际上这种阻碍关系正是能量的转化和守恒定律在电磁感应现象中的具体体现，故 D 正确。
}

\item
%% number: 15
%% paperName: 2019年全国理综Ⅲ第 15 题 6 分
%% typeId: 1
%% type: 单选题
%% chapter: 曲线运动
%% point: 天体运动
%% method: 无
%% score: 6
%% degree: 750
%% duplicateId: 0
%% body:
金星、地球和火星绕太阳的公转均可视为匀速圆周运动，它们的向心加速度大小分别为 $a_{\text{金}}$、$a_{\text{地}}$、$a_{\text{火}}$，它们沿轨道运行的速率分别为 $v_{\text{金}}$、$v_{\text{地}}$、$v_{\text{火}}$。已知它们的轨道半径 $R_{\text{金}} < R_{\text{地}} < R_{\text{火}}$，由此可以判定\xzanswer{A}
\fourchoices[answer=A]
{$a_{\text{金}} > a_{\text{地}} > a_{\text{火}}$}
{$a_{\text{火}} > a_{\text{地}} > a_{\text{金}}$}
{$v_{\text{地}} > v_{\text{火}} > v_{\text{金}}$}
{$v_{\text{火}} > v_{\text{地}} > v_{\text{金}}$}

%% answer: A
%% memo:
\memoanswer{
万有引力提供行星绕太阳做匀速圆周运动的向心力，由匀速圆周运动规律及牛顿第二定律得 $G\frac{Mm}{R^{2}}=m\frac{v^{2}}{R}=ma$，解得 $v=\sqrt{\frac{GM}{R}}$，$a=\frac{GM}{R^{2}}$。由于 $R_{\text{金}}<R_{\text{地}}<R_{\text{火}}$，所以 $v_{\text{金}}>v_{\text{地}}>v_{\text{火}}$，$a_{\text{金}}>a_{\text{地}}>a_{\text{火}}$，故 A 正确。
}

\item
%% number: 16
%% paperName: 2019年全国理综Ⅲ第 16 题 6 分
%% typeId: 1
%% type: 单选题
%% chapter: 力物体的平衡
%% point: 三力平衡
%% method: 无
%% score: 6
%% degree: 550
%% duplicateId: 0
%% body:
用卡车运输质量为 $m$ 的匀质圆筒状工件，为使工件保持固定，将其置于两光滑斜面之间，如图所示。两斜面 I、Ⅱ 固定在车上，倾角分别为 30${^\circ}$ 和 60${^\circ}$。重力加速度为 $g$。当卡车沿平直公路匀速行驶时，圆筒对斜面 I、Ⅱ 压力的大小分别为 $F_{1}$、$F_{2}$，则\xzanswer{D}
\onepicture[w=6.0cm]{figs/16.png}
\fourchoices[answer=D]
{$F_{1} = \frac{{\sqrt 3 }}{3}mg$，$F_{2} = \frac{{\sqrt 3 }}{2}mg$}
{$F_{1} = \frac{{\sqrt 3 }}{2}mg$，$F_{2} = \frac{{\sqrt 3 }}{3}mg$}
{$F_{1} = \frac{1}{2}mg$，$F_{2} = \frac{{\sqrt 3 }}{2}mg$}
{$F_{1} = \frac{{\sqrt 3 }}{2}mg$，$F_{2} = \frac{1}{2}mg$}

%% answer: D
%% memo:
\memoanswer{
以匀质圆筒为研究对象，把重力分解为压紧斜面Ⅰ的压力 $F_{1}$ 和压紧斜面Ⅱ的压力 $F_{2}$，如图所示，则有 $F_{1}=mg\cos30^{\circ}=\frac{\sqrt{3}}{2}mg$，$F_{2}=mg\sin30^{\circ}=\frac{1}{2}mg$，故 D 正确。

\onepicture[w=4.5cm]{figs/16_an.png}
}

\item
%% number: 17
%% paperName: 2019年全国理综Ⅲ第 17 题 6 分
%% typeId: 1
%% type: 单选题
%% chapter: 功和能
%% point: 功能中的图象
%% method: 无
%% score: 6
%% degree: 450
%% duplicateId: 0
%% body:
从地面竖直向上抛出一物体，物体在运动过程中除受到重力外，还受到一大小不变、方向始终与运动方向相反的外力作用。距地面高度 $h$ 在 $3\Um$ 以内时，物体上升、下落过程中动能 $E_{k}$ 随 $h$ 的变化如图所示。重力加速度取 $10\Umsq$。该物体的质量为\xzanswer{C}
\onepicture[w=6.5cm]{\includesvg{figs/17}}
\fourchoices[answer=C]
{$2\Ukg$}
{$1.5\Ukg$}
{$1\Ukg$}
{$0.5\Ukg$}

%% answer: C
%% memo:
\memoanswer{
由动能定理得 $W_{\text{合}}=\Delta E_{k}$，物体在上升过程中有 $-(mg+F)\times3\Um=36\UJ-72\UJ$，在下落过程中有 $(mg-F)\times3\Um=48\UJ-24\UJ$，联立解得该物体的质量 $m=1\Ukg$，故 C 正确。
}

\item
%% number: 18
%% paperName: 2019年全国理综Ⅲ第 18 题 6 分
%% typeId: 1
%% type: 单选题
%% chapter: 磁场
%% point: 带电粒子在磁场中的运动
%% method: 无
%% score: 6
%% degree: 550
%% duplicateId: 0
%% body:
如图，在坐标系的第一和第二象限内存在磁感应强度大小分别为 $\frac{B}{2}$ 和 $B$、方向均垂直于纸面向外的匀强磁场。一质量为 $m$、电荷量为 $q$（$q$ > 0）的粒子垂直于 $x$ 轴射入第二象限，随后垂直于 $y$ 轴进入第一象限，最后经过 $x$ 轴离开第一象限。粒子在磁场中运动的时间为\xzanswer{B}
\onepicture[w=6.5cm]{figs/18.png}
\fourchoices[answer=B]
{$\frac{{5\pi m}}{{6qB}}$}
{$\frac{{7\pi m}}{{6qB}}$}
{$\frac{{11\pi m}}{{6qB}}$}
{$\frac{{13\pi m}}{{6qB}}$}

%% answer: B
%% memo:
\memoanswer{
在第二象限中，带电粒子在匀强磁场中做匀速圆周运动，洛伦兹力提供向心力，由于粒子垂直于 $x$ 轴射入第二象限，又垂直于 $y$ 轴进入第一象限，如图所示，粒子在第二象限磁场中经过 $\frac{1}{4}$ 圆周，原点 O 即为圆心，设粒子的入射速度为 $v$，则有 $qvB=m\frac{v^2}{R}$，解得粒子运动的轨迹半径 $R=\frac{mv}{qB}$。由于第一象限内的匀强磁场的磁感应强度为 $\frac{1}{2}B$，则粒子运动的轨迹半径 $R'=2R$，圆心一定在 $y$ 轴上，画出粒子的运动轨迹，如图所示。由几何关系知 $\cos\theta=\frac{R}{2R}=\frac{1}{2}$，解得 $\theta=60^\circ$。因为粒子在磁场中运动周期 $T=\frac{2\pi m}{qB}$，所以粒子在磁场中运动的总时间为 $t=\frac{1}{4}T_1+\frac{60^\circ}{360^\circ}T_2=\frac{1}{4}\times\frac{2\pi m}{qB}+\frac{1}{6}\times\frac{2\pi m}{\frac{1}{2}qB}=\frac{7\pi m}{6qB}$，故 B 正确。

\onepicture[w=5.5cm]{figs/18_an.png}
}

\item
%% number: 19
%% paperName: 2019年全国理综Ⅲ第 19 题 6 分
%% typeId: 2
%% type: 多选题
%% chapter: 电磁感应
%% point: 电磁感应中图象
%% method: 无
%% score: 6
%% degree: 450
%% duplicateId: 0
%% body:
如图，方向竖直向下的匀强磁场中有两根位于同一水平面内的足够长的平行金属导轨，两相同的光滑导体棒 ab、cd 静止在导轨上。$t$ = 0 时，棒 ab 以初速度 $v_{0}$ 向右滑动。运动过程中，ab、cd 始终与导轨垂直并接触良好，两者速度分别用 $v_{1}$、$v_{2}$ 表示，回路中的电流用 $I$ 表示。下列图像中可能正确的是\xzanswer{AC}
\onepicture[w=6.0cm]{\includesvg{figs/19a}}
\fourchoices[answer=AC, ispicture=true, h=2.5cm]
{figs/19b.png}
{figs/19c.png}
{figs/19d.png}
{figs/19e.png}

%% answer: AC
%% memo:
\memoanswer{
以两导体棒为研究对象，在导体棒运动过程中，两导体棒所受的安培力大小相等，方向相反，且不受其他水平外力作用，在水平方向两导体棒组成的系统动量守恒，对系统有 $mv_{0}=2mv$，解得两导体棒运动的末速度为 $v=\frac{1}{2}v_{0}$，棒 ab 做变减速运动，棒 cd 做变加速运动，稳定时两导体棒的加速度均为零，一起向右做匀速运动，故 A 正确 B 错误；ab 棒和 cd 棒最后做匀速运动，棒与导轨组成的回路磁通量不变化，不会产生感应电流，故 C 正确 D 错误。
}

\item
%% number: 20
%% paperName: 2019年全国理综Ⅲ第 20 题 6 分
%% typeId: 2
%% type: 多选题
%% chapter: 运动和力
%% point: 牛顿定律中的图象
%% method: 无
%% score: 6
%% degree: 450
%% duplicateId: 0
%% body:
如图（a），物块和木板叠放在实验台上，物块用一不可伸长的细绳与固定在实验台上的力传感器相连，细绳水平。$t$ = 0 时，木板开始受到水平外力 $F$ 的作用，在 $t=4\Us$ 时撤去外力。细绳对物块的拉力 $f$ 随时间 $t$ 变化的关系如图（b）所示，木板的速度 $v$ 与时间 $t$ 的关系如图（c）所示。木板与实验台之间的摩擦可以忽略。重力加速度取 $g=10\Umsq$。由题给数据可以得出\xzanswer{AB}
\threepicture[labelA=2019全国理综Ⅲ20a, labelB=2019全国理综Ⅲ20b, labelC=2019全国理综Ⅲ20c, wA=6.0cm, wB=4.6cm, wC=4.6cm]
{figs/20a.png}
{figs/20b.png}
{figs/20c.png}
\fourchoices[answer=AB]
{木板的质量为 $1\Ukg$}
{$2\Us\sim4\Us$ 内，力 $F$ 的大小为 $0.4\UN$}
{$0\sim2\Us$ 内，力 $F$ 的大小保持不变}
{物块与木板之间的动摩擦因数为 0.2}

%% answer: AB
%% memo:
\memoanswer{
设木板质量为 $M$，物块质量为 $m$，在 $0\sim2\Us$ 内，结合图(b)(c)可知在 $F$ 作用下，木板与物块之间有静摩擦力作用，细绳对物块的拉力与静摩擦力平衡，可知力传感器测量的 $f$ 的大小等于摩擦力的大小，随着 $f$ 逐渐增大，木板保持静止状态，对木板和物块整体分析，可知力 $F$ 逐渐增大，当力 $F$ 大于木板与物块间的最大静摩擦力时，两者开始相对滑动，故 C 错误；$2\Us$ 后木板与物块之间有滑动摩擦力，大小为 $\mu mg=0.2\UN$，在 $2\sim4\Us$ 内，以木板为研究对象，利用牛顿第二定律有 $F-\mu mg=Ma'$，由图(c)可知 $a'=\frac{0.4-0}{4-2}\Umsq=0.2\Umsq$，在 $4\sim5\Us$ 内木板只在滑动摩擦力作用下做匀减速运动，根据牛顿第二定律有 $\mu mg=Ma$，由图(c)可知 $a=\left|\frac{0.2-0.4}{5-4}\right|\Umsq=0.2\Umsq$，联立解得 $F=0.4\UN$，$M=1\Ukg$，由于物块质量不确定，所以动摩擦因数的大小不确定，故 AB 正确 D 错误。
}

\item
%% number: 21
%% paperName: 2019年全国理综Ⅲ第 21 题 6 分
%% typeId: 2
%% type: 多选题
%% chapter: 电场
%% point: 电场线
%% method: 无
%% score: 6
%% degree: 450
%% duplicateId: 0
%% body:
如图，电荷量分别为 $q$ 和 $-q$（$q$ > 0）的点电荷固定在正方体的两个顶点上，a、b 是正方体的另外两个顶点。则\xzanswer{BC}
\onepicture[w=6.0cm]{\includesvg{figs/21}}
\fourchoices[answer=BC]
{a 点和 b 点的电势相等}
{a 点和 b 点的电场强度大小相等}
{a 点和 b 点的电场强度方向相同}
{将负电荷从 a 点移到 b 点，电势能增加}

%% answer: BC
%% memo:
\memoanswer{
正点电荷的场强方向沿径向向外，负点电荷的场强方向沿径向向里，故 $q$ 在 b 点的场强方向沿 $q$ 指向 b，大小为 $\frac{kq}{r^2}$，$-q$ 在 b 点的场强方向沿 b 指向 $-q$，大小为 $\frac{kq}{2r^2}$，两个场强方向垂直，同理得 $q$ 在 a 点的场强与 $-q$ 在 b 点的场强相同，$-q$ 在 a 点的场强与 $q$ 在 b 点的场强相同，故 a、b 两点的场强大小、方向都相同，故 BC 正确；点电荷的等势面为以点电荷为球心的球面，故 a 点的电势与 b 点正上方顶点的电势相同，将负电荷从 b 点正上方的顶点沿边线移动到 b 点，负电荷所受的两个电场力都做正功，故负电荷的电势能降低，a 点电势低于 b 点，故 AD 错误。

\onepicture[w=4.6cm]{figs/21_an1.png}
\onepicture[w=3.8cm]{figs/21_an2.png}
}

\item
%% number: 22
%% paperName: 2019年全国理综Ⅲ第 22 题 5 分
%% typeId: 4
%% type: 实验题
%% chapter: 直线运动
%% point: 打点计时器公式
%% method: 无
%% score: 5
%% degree: 750
%% duplicateId: 0
%% body:
甲乙两位同学设计了利用数码相机的连拍功能测重力加速度的实验。实验中，甲同学负责释放金属小球，乙同学负责在小球自由下落的时候拍照。已知相机每间隔 $0.1\Us$ 拍 1 幅照片。
\begin{enumerate}
	\item
	若要从拍得的照片中获取必要的信息，在此实验中还必须使用的器材是\tkanswer[1cm]{A}。（填正确答案标号）
	\fourchoices[answer=A]
	{米尺}
	{秒表}
	{光电门}
	{天平}
	\item
	简述你选择的器材在本实验中的使用方法。
	\item
	实验中两同学由连续 3 幅照片上小球的位置 a、b 和 c 得到 ab = $24.5\Ucm$、ac = $58.7\Ucm$，则该地的重力加速度大小为 $g=$\tkanswer[1.5cm]{9.7} $\Umsq$。（保留 2 位有效数字）
\end{enumerate}

%% answer:
\jdanswer{
\begin{enumerate}
	\item
	A

	\item
	将米尺竖直放置，使小球下落时尽量靠近米尺

	\item
	9.7

\end{enumerate}
}
%% memo:
\memoanswer{
（1）利用匀变速直线运动规律测量重力加速度，由于时间已知，根据 $h=\frac{1}{2}gt^{2}$，可知还需要测量长度，故为了从照片中获取必要的信息，还必须使用的器材是米尺，故 A 正确。

（2）将米尺竖直放置，使小球下落时尽量靠近米尺，用米尺测量出照片上相邻小球间的距离。

（3）因 $ab=24.5\Ucm$、$ac=58.7\Ucm$，所以 $bc=ac-ab=34.2\Ucm$，由匀变速直线运动的规律 $\Delta x=gT^2$ 可得 $g=\frac{bc-ab}{T^2}$，代入数据得 $g=\frac{34.2-24.5}{0.1^2}\times10^{-2}\Umsq=9.7\Umsq$。
}

\item
%% number: 23
%% paperName: 2019年全国理综Ⅲ第 23 题 10 分
%% typeId: 4
%% type: 实验题
%% chapter: 电路
%% point: 多用表的使用
%% method: 无
%% score: 10
%% degree: 450
%% duplicateId: 0
%% body:
某同学欲将内阻为 $98.5\UO$、量程为 $100\UuA$ 的电流表改装成欧姆表并进行刻度和校准，要求改装后欧姆表的 $15\UkO$ 刻度正好对应电流表表盘的 $50\UuA$ 刻度。可选用的器材还有：定值电阻 $R_{0}$（阻值 $14\UkO$），滑动变阻器 $R_{1}$（最大阻值 $1500\UO$），滑动变阻器 $R_{2}$（最大阻值 $500\UO$），电阻箱（$0\sim99999.9\UO$），干电池（$E=1.5\UV$，$r=1.5\UO$），红、黑表笔和导线若干。
\begin{enumerate}
	\item
	欧姆表设计

	将图（a）中的实物连线组成欧姆表。欧姆表改装好后，滑动变阻器 $R$ 接入电路的电阻应为\tkanswer[1.5cm]{900}$\UO$；滑动变阻器选\tkanswer[1cm]{$R_{1}$}（填“$R_{1}$”或“$R_{2}$”）。
	\item
	刻度欧姆表表盘

	通过计算，对整个表盘进行电阻刻度，如图（b）所示。表盘上 a、b 处的电流刻度分别为 25 和 75，则 a、b 处的电阻刻度分别为\tkanswer[1cm]{45}、\tkanswer[1cm]{5}。
	\item
	校准

	红、黑表笔短接，调节滑动变阻器，使欧姆表指针指向\tkanswer[1cm]{0}$\UkO$ 处；将红、黑表笔与电阻箱连接，记录多组电阻箱接入电路的电阻值及欧姆表上对应的测量值，完成校准数据测量。若校准某刻度时，电阻箱旋钮位置如图（c）所示，则电阻箱接入的阻值为\tkanswer[2cm]{35000.0}$\UO$。
\end{enumerate}
\threepicture[labelA=2019全国理综Ⅲ23a, labelB=2019全国理综Ⅲ23b, labelC=2019全国理综Ⅲ23c, wA=4.0cm, wB=5.5cm, wC=4.0cm, align=right]
{figs/23a.png}
{figs/23b.png}
{figs/23c.png}

%% answer:
\jdanswer{
\begin{enumerate}
	\item
	见解析

	\item
	45，5

	\item
	0，35000.0

\end{enumerate}
}
%% memo:
\memoanswer{
（1）实物连线如图所示。欧姆表内部由电源、定值电阻、电流表和滑动变阻器构成，中值电阻为 $15\UkO$，即欧姆调零时有 $I_{A}=\frac{E}{R_{\text{内}}}$，半偏时有 $0.5I_{A}=\frac{E}{R_{\text{内}}+R_{\text{中}}}$，解得欧姆表内阻为 $R_{\text{内}}=15\UkO$，滑动变阻器接入电路的阻值 $R=R_{\text{内}}-R_{0}-R_{A}-r=900\UO$，则滑动变阻器应选 $R_{1}$。

\onepicture[w=5.5cm]{figs/23_an.png}

（2）由闭合电路欧姆定律得，指针指 a 处时电流为 $I_{a}=\frac{E}{R_{\text{内}}+R_{a}}=0.25I_{A}$，指 b 处时电流为 $I_{b}=\frac{E}{R_{\text{内}}+R_{b}}=0.75I_{A}$，又 $I_{A}=\frac{E}{R_{\text{内}}}$，联立解得 $R_{a}=45\UkO$，$R_{b}=5\UkO$，故 a、b 处的电阻刻度分别为 45、5。

（3）欧姆调零时，应该使电流表指针指向电流满偏处，此时外电路的电阻为零，欧姆表指针指向 $0\UkO$ 处，电阻箱的读数为 $3\times10\UkO+5\times1\UkO+0\times100\UO+0\times10\UO+0\times1\UO+0\times0.1\UO=35000.0\UO$。
}

\item
%% number: 24
%% paperName: 2019年全国理综Ⅲ第 24 题 12 分
%% typeId: 6
%% type: 计算题
%% chapter: 电场
%% point: 带电粒子的偏转
%% method: 无
%% score: 12
%% degree: 550
%% duplicateId: 0
%% body:
空间存在一方向竖直向下的匀强电场，O、P 是电场中的两点。从 O 点沿水平方向以不同速度先后发射两个质量均为 $m$ 的小球 A、B。A 不带电，B 的电荷量为 $q$（$q$ > 0）。A 从 O 点发射时的速度大小为 $v_{0}$，到达 P 点所用时间为 $t$；B 从 O 点到达 P 点所用时间为 $\frac{t}{2}$。重力加速度为 $g$，求：
\begin{enumerate}
	\item
	电场强度的大小；
	\item
	B 运动到 P 点时的动能。
\end{enumerate}

%% answer:
\jdanswer{
\begin{enumerate}
	\item
	$E = \frac{{3mg}}{q}$

	\item
	$E_{k}$ = 2$m$($v_{0}^{2} + g^{2}t^{2}$)

\end{enumerate}
}
%% memo:
\memoanswer{
（1）设电场强度的大小为 $E$，小球 B 运动的加速度为 $a$。根据牛顿定律、运动学公式和题给条件，有 $mg+qE=ma$①，$\frac{1}{2}a\left(\frac{t}{2}\right)^2=\frac{1}{2}gt^2$②，解得 $E=\frac{3mg}{q}$③。

（2）设 B 从 O 点发射时的速度为 $v_{1}$，到达 P 点时的动能为 $E_{k}$，O、P 两点的高度差为 $h$，根据动能定理有 $E_{k}-\frac{1}{2}mv_{1}^{2}=mgh+qEh$④，且有 $v_{1}\frac{t}{2}=v_{0}t$⑤，$h=\frac{1}{2}gt^{2}$⑥，联立③④⑤⑥式得 $E_{k}=2m(v_{0}^{2}+g^{2}t^{2})$⑦。
}

\item
%% number: 25
%% paperName: 2019年全国理综Ⅲ第 25 题 20 分
%% typeId: 6
%% type: 计算题
%% chapter: 运动和力
%% point: 动量和能量
%% method: 无
%% score: 20
%% degree: 350
%% duplicateId: 0
%% body:
静止在水平地面上的两小物块 A、B，质量分别为 $m_{A}=1.0\Ukg$，$m_{B}=4.0\Ukg$；两者之间有一被压缩的微型弹簧，A 与其右侧的竖直墙壁距离 $l=1.0\Um$，如图所示。某时刻，将压缩的微型弹簧释放，使 A、B 瞬间分离，两物块获得的动能之和为 $E_{k}=10.0\UJ$。释放后，A 沿着与墙壁垂直的方向向右运动。A、B 与地面之间的动摩擦因数均为 $\mu=0.20$。重力加速度取 $g=10\Umsq$。A、B 运动过程中所涉及的碰撞均为弹性碰撞且碰撞时间极短。
\begin{enumerate}
	\item
	求弹簧释放后瞬间 A、B 速度的大小；
	\item
	物块 A、B 中的哪一个先停止？该物块刚停止时 A 与 B 之间的距离是多少？
	\item
	A 和 B 都停止后，A 与 B 之间的距离是多少？
\end{enumerate}
\onepicture[w=8cm, align=right]{\includesvg{figs/25}}

%% answer:
\jdanswer{
\begin{enumerate}
	\item
	$v_{A}=4.0\Ums$，$v_{B}=1.0\Ums$

	\item
	B 先停止，$0.50\Um$

	\item
	$0.91\Um$

\end{enumerate}
}
%% memo:
\memoanswer{
（1）设弹簧释放瞬间 A 和 B 的速度大小分别为 $v_{A}$、$v_{B}$，以向右为正，由动量守恒定律和题给条件得 $0=m_{A}v_{A}-m_{B}v_{B}$①，$E_{k}=\frac{1}{2}m_{A}v_{A}^{2}+\frac{1}{2}m_{B}v_{B}^{2}$②，联立①②式并代入题给数据得 $v_{A}=4.0\Ums$、$v_{B}=1.0\Ums$③。

（2）A、B 两物块与地面间的动摩擦因数相等，因而两者滑动时加速度大小相等，设为 $a$。假设 A 和 B 发生碰撞前，已经有一个物块停止，此物块应为弹簧释放后速度较小的 B。设从弹簧释放到 B 停止所需时间为 $t$，B 向左运动的路程为 $s_{B}$，则有 $m_{B}a=\mu m_{B}g$④，$s_{B}=v_{B}t-\frac{1}{2}at^{2}$⑤，$v_{B}-at=0$⑥。在时间 $t$ 内，A 可能与墙发生弹性碰撞，碰撞后 A 将向左运动，碰撞并不改变 A 的速度大小，所以无论此碰撞是否发生，A 在时间 $t$ 内的路程 $s_{A}$ 都可以表示为 $s_{A}=v_{A}t-\frac{1}{2}at^{2}$⑦，联立③④⑤⑥⑦式并代入题给数据得 $s_{A}=1.75\Um$、$s_{B}=0.25\Um$⑧，这表明在时间 $t$ 内 A 已与墙壁发生碰撞，但没有与 B 发生碰撞，此时 A 位于出发点右边 $0.25\Um$ 处，B 位于出发点左边 $0.25\Um$ 处，两物块之间距离 $s$ 为 $s=0.25\Um+0.25\Um=0.50\Um$⑨。

（3）$t$ 时刻后 A 将继续向左运动，假设它能与静止的 B 碰撞，碰撞时速度的大小为 $v_{A}'$，由动能定理得 $\frac{1}{2}m_{A}v_{A}'^{2}-\frac{1}{2}m_{A}v_{A}^{2}=-\mu m_{A}g(2l+s_{B})$⑩，联立③⑧⑩式并代入题给数据得 $v_{A}'=\sqrt{7}\Ums$\ccnum{⑪}，故 A 与 B 将发生碰撞。设碰撞后 A、B 的速度分别为 $v_{A}''$ 和 $v_{B}''$，由动量守恒定律与机械能守恒定律有 $m_{A}(-v_{A}')=m_{A}v_{A}''+m_{B}v_{B}''$\ccnum{⑫}，$\frac{1}{2}m_{A}v_{A}'^{2}=\frac{1}{2}m_{A}v_{A}''^{2}+\frac{1}{2}m_{B}v_{B}''^{2}$\ccnum{⑬}，联立\ccnum{⑪}\ccnum{⑫}\ccnum{⑬}式并代入题给数据得 $v_{A}''=\frac{3\sqrt{7}}{5}\Ums$、$v_{B}''=-\frac{2\sqrt{7}}{5}\Ums$\ccnum{⑭}，这表明碰撞后 A 将向右运动，B 继续向左运动。设碰撞后 A 向右运动距离为 $s_{A}'$ 时停止，B 向左运动距离为 $s_{B}'$ 时停止，由运动学公式得 $2as_{A}'=v_{A}''^{2}$、$2as_{B}'=v_{B}''^{2}$\ccnum{⑮}，由④\ccnum{⑭}\ccnum{⑮}式及题给数据得 $s_{A}'=0.63\Um$、$s_{B}'=0.28\Um$\ccnum{⑯}，$s_{A}'$ 小于碰撞处到墙壁的距离。由上式可得两物块停止后的距离 $s'=s_{A}'+s_{B}'=0.91\Um$\ccnum{⑰}。
}

\item
%% number: 33
%% paperName: 2019年全国理综Ⅲ第 33 题 10 分
%% typeId: 6
%% type: 计算题
%% chapter: 气体性质
%% point: 气体多过程
%% method: 无
%% score: 10
%% degree: 550
%% duplicateId: 0
%% body:
\begin{enumerate}
	\item
	（5 分）用油膜法估算分子大小的实验中，首先需将纯油酸稀释成一定浓度的油酸酒精溶液，稀释的目的是\tkanswer[3cm]{使油酸在浅盘的水面上容易形成一块单分子层油膜}。实验中为了测量出一滴已知浓度的油酸酒精溶液中纯油酸的体积，可以\tkanswer[3cm]{把油酸酒精溶液一滴一滴地滴入小量筒中，测出 $1\UmL$ 油酸酒精溶液的滴数，得到一滴溶液中纯油酸的体积}。为得到油酸分子的直径，还需测量的物理量是\tkanswer[2.5cm]{单分子层油膜的面积}。
	\item
	（10 分）如图，一粗细均匀的细管开口向上竖直放置，管内有一段高度为 $2.0\Ucm$ 的水银柱，水银柱下密封了一定量的理想气体，水银柱上表面到管口的距离为 $2.0\Ucm$。若将细管倒置，水银柱下表面恰好位于管口处，且无水银滴落，管内气体温度与环境温度相同。已知大气压强为 $76\UcmHg$，环境温度为 $296\UK$。
	\begin{enumerate}
		\item
		求细管的长度；
		\item
		若在倒置前，缓慢加热管内被密封的气体，直到水银柱的上表面恰好与管口平齐为止，求此时密封气体的温度。
	\end{enumerate}
\end{enumerate}
\onepicture[w=3cm, align=right]{\includesvg{figs/33}}

%% answer:
\jdanswer{
\begin{enumerate}
	\item
	使油酸在浅盘的水面上容易形成一块单分子层油膜；把油酸酒精溶液一滴一滴地滴入小量筒中，测出 $1\UmL$ 油酸酒精溶液的滴数，得到一滴溶液中纯油酸的体积；单分子层油膜的面积

	\item
	(i) $41\Ucm$；(ii) $312\UK$

\end{enumerate}
}
%% memo:
\memoanswer{
（1）纯油酸的粘稠度比较大，在水面上不容易散开，稀释的目的是尽量降低油酸的浓度，使滴在水面上的油酸尽量散开，形成单分子层油膜，同时酒精易挥发，不影响测量结果。实验中为了测量出一滴已知浓度的油酸酒精溶液中纯油酸的体积，可以把一定浓度的油酸酒精溶液一滴一滴地滴入小量筒，记下滴入溶液滴数、量筒内油酸酒精溶液的体积，则可以计算出一滴溶液中纯油酸的体积。为得到油酸分子的直径，还需测量的物理量是单分子层油膜的面积。

（2）(i) 设细管的长度为 $L$，横截面的面积为 $S$，水银柱高度为 $h$；初始时，设水银柱上表面到管口的距离为 $h_{1}$，被密封气体的体积为 $V$，压强为 $p$；细管倒置时，气体体积为 $V_{1}$，压强为 $p_{1}$。由玻意耳定律得 $pV=p_{1}V_{1}$①，由力的平衡条件得 $p=p_{0}+\rho gh$②，$p_{1}=p_{0}-\rho gh$③，式中 $\rho$、$g$ 分别为水银的密度和重力加速度的大小，$p_{0}$ 为大气压强。由题意得 $V=S(L-h_{1}-h)$④，$V_{1}=S(L-h)$⑤，由①②③④⑤式和题给条件得 $L=41\Ucm$⑥。

\onepicture[w=3cm]{\includesvg{figs/33_an}}

(ii) 设气体被加热前后的温度分别为 $T_{0}$ 和 $T$，由盖-吕萨克定律得 $\frac{V}{T_{0}}=\frac{V_{1}}{T}$⑦，由④⑤⑥⑦式和题给数据得 $T=312\UK$⑧。
}

\item
%% number: 34
%% paperName: 2019年全国理综Ⅲ第 34 题 10 分
%% typeId: 6
%% type: 计算题
%% chapter: 光学
%% point: 光的折射
%% method: 无
%% score: 10
%% degree: 550
%% duplicateId: 0
%% body:
\begin{enumerate}
	\item
	（5 分）水槽中，与水面接触的两根相同细杆固定在同一个振动片上。振动片做简谐振动时，两根细杆周期性触动水面形成两个波源。两波源发出的波在水面上相遇。在重叠区域发生干涉并形成了干涉图样。关于两列波重叠区域内水面上振动的质点，下列说法正确的是\xzanswer{BDE}
	\fivechoices[answer=BDE]
	{不同质点的振幅都相同}
	{不同质点振动的频率都相同}
	{不同质点振动的相位都相同}
	{不同质点振动的周期都与振动片的周期相同}
	{同一质点处，两列波的相位差不随时间变化}
	\item
	（10 分）如图，直角三角形 ABC 为一棱镜的横截面，$\angle$A = 90${^\circ}$，$\angle$B = 30${^\circ}$。一束光线平行于底边 BC 射到 AB 边上并进入棱镜，然后垂直于 AC 边射出。
	\begin{enumerate}
		\item
		求棱镜的折射率；
		\item
		保持 AB 边上的入射点不变，逐渐减小入射角，直到 BC 边上恰好有光线射出。求此时 AB 边上入射角的正弦。
	\end{enumerate}
\end{enumerate}
\onepicture[w=6cm, align=right]{\includesvg{figs/34}}

%% answer:
\jdanswer{
\begin{enumerate}
	\item
	BDE

	\item
	(i) $\sqrt 3 $；(ii) $\frac{\sqrt 3 - \sqrt 2 }{2}$

\end{enumerate}
}
%% memo:
\memoanswer{
（1）发生干涉时，有的质点振动加强，有的质点振动减弱，不同质点振幅不一定相同，故 A 错误；两列波叠加区域，介质中的质点同时参与这两列波引起的振动，质点的位移等于两列波单独传播时引起的位移的矢量和，两个波源的周期相同，振幅相同，故所有质点振动的周期相同，都等于波源的周期，频率 $f=\frac{1}{T}$ 也相同，故 BD 正确；不同质点的振动有的加强，有的减弱，振动的相位不一定相同，故 C 错误；两列波的频率相同，同一质点处，两列波的波程差恒定，所以相位差不变，故 E 正确。

（2）(i) 光路图及相关量如图所示。光束在 AB 边上折射，由折射定律得 $\frac{\sin i}{\sin\alpha}=n$①，式中 $n$ 是棱镜的折射率。由几何关系得 $\alpha+\beta=60^{\circ}$②，由几何关系和反射定律得 $\beta=\beta'=\angle B$③，联立①②③式，并代入 $i=60^{\circ}$ 得 $n=\sqrt{3}$④。

\onepicture[w=8cm]{figs/34_an.png}

(ii) 设改变后的入射角为 $i'$，折射角为 $\alpha'$，由折射定律得 $\frac{\sin i'}{\sin\alpha'}=n$⑤，依题意，光束在 BC 边上的入射角为全反射的临界角 $\theta_{C}$，且 $\sin\theta_{C}=\frac{1}{n}$⑥，由几何关系得 $\theta_{C}=\alpha'+30^{\circ}$⑦，由④⑤⑥⑦式得入射角的正弦为 $\sin i'=\frac{\sqrt{3}-\sqrt{2}}{2}$⑧。
}

\end{enumerate}

\end{document}
